% 第4章 键盘输入和屏幕输出（教材第4章 4.1--4.6）
\lstset{language=C}

\section{单个字符的输入/输出}

\subsection{字符常量}

字符常量是用一对单引号括起来的单个字符。C 语言中，字符在内存里存的是它的 \textbf{ASCII 码值}，因此字符与整数在数值上是等价的，可以直接参与整型运算：

\begin{lstlisting}
char ch = 'A';
printf("%c %d\n", ch, ch);      /* A 65 */
ch = ch + 32;                    /* 'A' -> 'a'，65 + 32 = 97 */
printf("%c %d\n", ch, ch);      /* a 97 */
\end{lstlisting}

观察 ASCII 码表可以发现一条常用规律：\textbf{小写字母的 ASCII 码值比相应大写字母大 32}，即 \texttt{'a'} 与 \texttt{'A'}、\texttt{'b'} 与 \texttt{'B'} 都相差 32。据此可轻便地实现大小写转换；反过来，\texttt{ch - 'a' + 'A'} 可把小写转成大写。这类“字符就是整数”的性质在字符分类、字符串处理中会反复用到。

字符常量书写有三条限制：单引号内\textbf{只能有一个字符}（\texttt{'ab'} 非法）；转义字符算一个字符（\texttt{'\textbackslash n'} 合法）；\texttt{'5'} 是字符（ASCII 值 53），而 \texttt{5} 是整数，两者不能混用。

\subsection{字符的输入/输出}

只处理单个字符时，标准库提供了一对专用函数：

\begin{lstlisting}
int getchar(void);        /* 从键盘读入一个字符，返回其 ASCII 值 */
int putchar(int c);       /* 向屏幕输出一个字符 */
\end{lstlisting}

\begin{lstlisting}
ch = getchar();     /* 读入一个字符，赋给 ch */
putchar(ch);        /* 输出该字符 */
putchar('\n');      /* 输出一个换行控制符 */
\end{lstlisting}

要点如下：

\begin{itemize}
  \item \texttt{getchar()} \textbf{没有参数}，它的返回值是读到的字符的 ASCII 码，所以既可以赋给 \texttt{char} 变量，也可以赋给 \texttt{int} 变量；
  \item \texttt{putchar()} 的参数是一个字符（或它的 ASCII 码值），一次只能输出一个字符；
  \item \texttt{getchar()} 从\textbf{输入缓冲区}读取字符，而不是每次按键都立即返回。用户按回车后，整行字符（\textbf{包括回车符}）才一起进入缓冲区；之后每次调用 \texttt{getchar()} 都直接从缓冲区取字符，只有缓冲区被读空后才会再次等待按键。这条机制是 4.4 节一系列问题的根源；
  \item 用 \texttt{printf("\%c", ch)} 与 \texttt{putchar(ch)} 输出一个字符的效果等价；而 \texttt{printf("\textbackslash n")} 与 \texttt{putchar('\textbackslash n')} 也等价，区别只在于前者写在双引号内（字符串），后者写在单引号内（单个字符）。
\end{itemize}

\section{数据的格式化屏幕输出}

\texttt{printf()} 的一般调用形式为
\[
\texttt{printf(格式控制字符串,\ 输出值参数表);}
\]

\textbf{格式控制字符串}是用双引号括起来的字符串，包含两类内容：\textbf{格式转换说明符}（以 \texttt{\%} 开头，指定数据的输出格式）和\textbf{普通字符}（原样输出）。\textbf{输出值参数表}是需要输出的数据项列表，可以是变量或表达式，之间用逗号分隔，其类型应与格式转换说明符\textbf{一一对应}。

\begin{table}[htbp]
\centering
\caption{printf 常用格式转换说明符}
\small
\begin{tabularx}{\textwidth}{@{}llX@{}}
\toprule
格式符 & 对应类型 & 说明 \\
\midrule
\texttt{\%d} & \texttt{int} & 十进制整数；\texttt{\%ld} 对应 \texttt{long} \\
\texttt{\%c} & \texttt{char} & 按字符输出；用 \texttt{\%d} 则输出其 ASCII 码值 \\
\texttt{\%f} & \texttt{float}/\texttt{double} & 小数形式，隐含输出 6 位小数 \\
\texttt{\%e} & 实型 & 指数形式 \\
\texttt{\%g} & 实型 & 自动选择较短形式，并去掉无意义的 0 \\
\texttt{\%s} & 字符串 & 输出字符串 \\
\texttt{\%o} / \texttt{\%x} & 整型 & 八进制 / 十六进制 \\
\texttt{\%u} & 无符号整型 & 十进制无符号 \\
\bottomrule
\end{tabularx}
\end{table}

\textbf{域宽与精度}可以写在 \texttt{\%} 与格式字符之间：\texttt{\%7.2f} 表示输出数据占 \textbf{7 个字符宽度}、保留 \textbf{2 位小数}。注意\textbf{小数点本身也占 1 个字符位置}，所以 \texttt{33.30} 共 5 个字符，前面会补 2 个空格。若数据实际宽度超过指定域宽，则按实际宽度输出，域宽不起作用。若精度写成 \texttt{.0}（如 \texttt{\%7.0f}），相当于输出实数的取整结果。

\begin{lstlisting}
double c = 33.300854, a = 88.247263;
printf("circumference = %f, area = %f\n", c, a);
/* 隐含 6 位小数：33.300854, 88.247263 */
printf("%7.2f, %7.2f\n", c, a);
/* 33.30,   88.25 —— 各占 7 列，右对齐，小数保留 2 位 */
printf("%-7.2f|\n", c);
/* 33.30  | —— 负号表示左对齐，右侧补空格 */
\end{lstlisting}

还有几个实用修饰：\texttt{-} 表示左对齐（默认右对齐），\texttt{+} 表示正数也显示正号，\texttt{0} 表示用 0 填充空位（如 \texttt{\%05d}）。\texttt{printf()} 的\textbf{返回值}是实际输出的字符个数，一般不用，但在拼接输出长度时有用。

\corepoint{\texttt{printf()} 最常见的两类错误：一是格式符与参数类型不匹配（用 \texttt{\%d} 输出实型、用 \texttt{\%f} 输出整型），输出结果是不可预期的乱值；二是格式符个数与参数个数不等，多出的参数被忽略、缺少的参数会读到栈上的垃圾值。写完一行 \texttt{printf}，从左到右数一遍 \texttt{\%} 的个数与逗号后面参数的个数是否相等。}

\section{数据的格式化键盘输入}

\texttt{scanf()} 的一般调用形式为
\[
\texttt{scanf(格式控制字符串,\ 参数地址表);}
\]

格式控制字符串同样由格式转换说明符和分隔符两部分组成；\textbf{参数地址表}给出的不是变量本身，而是变量\textbf{的地址}，因此每个参数前通常要加取地址运算符 \texttt{\&}。

\begin{lstlisting}
int a, b;
double x;
scanf("%d%d", &a, &b);      /* 输入两个整数，用空白分隔 */
scanf("%lf", &x);            /* 读入 double 必须用 %lf */
\end{lstlisting}

\begin{table}[htbp]
\centering
\caption{scanf 常用格式转换说明符及其输入规则}
\small
\begin{tabularx}{\textwidth}{@{}llX@{}}
\toprule
格式符 & 对应类型 & 输入规则 \\
\midrule
\texttt{\%d} & \texttt{int} & 跳过前导空白，读到非数字字符为止 \\
\texttt{\%f} & \texttt{float} & 读实数 \\
\texttt{\%lf} & \texttt{double} & \texttt{double} 变量必须用 \texttt{\%lf} \\
\texttt{\%c} & \texttt{char} & \textbf{不跳过空白}，连空格、回车都当有效字符读入 \\
\texttt{\%s} & 字符数组 & 跳过前导空白，读到空白字符为止 \\
\texttt{\%2d} & \texttt{int} & 限定域宽，自动按指定宽度截取输入数据 \\
\texttt{\%*d} & —— & 忽略符，读入该数据但不赋给任何变量 \\
\bottomrule
\end{tabularx}
\end{table}

\textbf{输入分隔规则}是最容易踩坑的地方：

\begin{itemize}
  \item 数值型（\texttt{\%d}、\texttt{\%f}）格式符之间用\textbf{空白字符分隔}时，输入数据之间可用空格、回车或制表符任一种分隔，不必与格式串严格一致；
  \item 若格式串中两个格式符之间写了\textbf{普通字符}（如逗号），则用户输入时必须\textbf{原样输入}该字符。\texttt{scanf("\%d,\%d", \&a, \&b)} 要求输入 \texttt{12,34}；
  \item 若格式串中写了提示文字（如 \texttt{"a=\%d,b=\%d"}），用户也必须把这些文字原样敲进去，这在实际使用中很不友好，通常改为先用 \texttt{printf} 提示、再只写格式符；
  \item \textbf{不要在格式串末尾加 \texttt{\textbackslash n}}。\texttt{scanf} 中的 \texttt{\textbackslash n} 不是“等待换行”，而是“跳过所有空白字符”，会导致程序一直等待输入；
  \item \texttt{scanf()} 的\textbf{返回值}是成功读入的数据个数，可用于判断输入是否合法（如 \texttt{if (scanf("\%d", \&n) != 1)}）。
\end{itemize}

\corepoint{\texttt{scanf} 忘记写 \texttt{\&} 是初学者最高频的致命错误。此时编译器把变量的\textbf{值}当作地址，数据被写到非法内存单元，真正的变量却没有被赋值，运行时可能直接崩溃（非法内存访问），也可能“看起来没报错但结果是垃圾值”。看到“输入了却没生效”，第一件事就是回去数 \texttt{\&}。}

\section{本章扩充内容}

\subsection{用 \%c 输入字符时存在的问题}

\texttt{\%c} 的特殊之处在于它\textbf{不跳过任何空白字符}。结合 4.1 节讲的缓冲区机制，就会出现这样的现象：

\begin{lstlisting}
int data1, data2;  char op;
printf("Please enter the expression data1+data2\n");
scanf("%d%c%d", &data1, &op, &data2);
printf("%d%c%d = %d\n", data1, op, data2, data1 + data2);
\end{lstlisting}

用户输入 \texttt{12 + 3} 时，期望 \texttt{op} 读到 \texttt{'+'}，实际却把空格读给了 \texttt{op}，于是 \texttt{data2} 读到的是 \texttt{'+'}（其 ASCII 值 43），输出变成 \texttt{12 3 = 3141} 之类的怪结果。

更常见的版本是连续读入不同类型的数据：

\begin{lstlisting}
scanf("%d", &a);      /* 用户输入 12 后按回车 */
scanf("%c", &b);      /* b 读到的是回车符，而不是下一个有效字符 */
\end{lstlisting}

原因同上一节：输入 \texttt{12} 后按下的\textbf{回车符留在了缓冲区}，下一次 \texttt{\%c} 调用直接从缓冲区取走它。这类错误的典型症状是“第二次输入被跳过了”。

\subsection{\%c 格式符存在问题的解决方法}

有两种标准做法，第二种更简单、可读性也更好：

\begin{lstlisting}
/* 方法 1：用 getchar() 把缓冲区中的回车符读掉 */
scanf("%d", &a);
getchar();               /* 吃掉上一次输入残留的换行符 */
scanf("%c", &b);

/* 方法 2：在 %c 前加一个空格，跳过前面所有空白字符 */
scanf("%d", &a);
scanf(" %c", &b);        /* 注意 %c 前面那个空格是关键 */
\end{lstlisting}

两种方法的差别在于：\texttt{getchar()} 只能吃掉\textbf{一个}字符，如果缓冲区里残留了多个空白（比如输入时多敲了几个空格），仍可能出错；而格式串中的空格表示“跳过任意多个空白字符”，更稳健。因此在混合读入数值与字符时，推荐统一写成 \texttt{" \%c"}。

\section{本章知识点小结}

\begin{itemize}
  \item 字符以 ASCII 码存储，与整数通用；大小写字母的码值相差 32；
  \item \texttt{getchar()}/\texttt{putchar()} 面向单个字符；\texttt{getchar()} 从输入缓冲区取字符，回车也进缓冲区；
  \item \texttt{printf} 由格式控制字符串与输出值参数表组成，格式符与参数类型、个数必须一一对应；域宽与精度写作 \texttt{\%m.nf}；
  \item \texttt{scanf} 的参数是\textbf{地址}，必须有 \texttt{\&}；数值型可用任意空白分隔，格式串中的普通字符必须原样输入；
  \item \texttt{\%c} 不跳过空白，混合输入时用“空格 + \texttt{\%c}”或 \texttt{getchar()} 清掉残留回车；
  \item \texttt{double} 的输入必须用 \texttt{\%lf}，而输出用 \texttt{\%f} 即可。
\end{itemize}

\section{本章常见错误小结}

\begin{table}[htbp]
\centering
\caption{第 4 章常见错误与纠正}
\small
\begin{tabularx}{\textwidth}{@{}lXl@{}}
\toprule
易错写法 & 错在哪里 & 纠正 \\
\midrule
\texttt{scanf("\%d", n);} & 漏写取地址符，数据写到非法地址 & \texttt{scanf("\%d", \&n);} \\
\texttt{scanf("\%d\textbackslash n", \&n);} & \texttt{\textbackslash n} 表示跳过所有空白，程序会一直等待输入 & 去掉 \texttt{\textbackslash n} \\
\texttt{scanf("\%f", \&x)} 而 \texttt{x} 是 \texttt{double} & 类型不匹配，读入错误 & 用 \texttt{\%lf} \\
\texttt{scanf("\%d,\%d", \&a,\&b)} 输入 \texttt{12 34} & 格式串中的逗号必须原样输入 & 输入 \texttt{12,34} 或改格式串 \\
连续 \texttt{scanf("\%d")} 后跟 \texttt{scanf("\%c")} & \texttt{\%c} 读到残留的回车符 & \texttt{scanf(" \%c", \&ch)} \\
\texttt{printf("\%d", 3.5)} & 格式符与参数类型不匹配 & 用 \texttt{\%f} 输出实型 \\
\texttt{printf("\%f", n)} 而 \texttt{n} 是 \texttt{int} & 同上，输出不可预期的值 & 换成 \texttt{\%d} \\
\texttt{putchar("a")} & 双引号是字符串，参数应为字符 & \texttt{putchar('a')} \\
\bottomrule
\end{tabularx}
\end{table}
